% Zone „Das kennst du schon“ – aus bank/hypothesentests/zone.jsonl (zusammenbau v0.1)
\einheitenkopf[zone][Kennst du schon]{Das kennst du schon}
\setcounter{aufgabe}{0}

%% TODO Titel: Ich-kann-Satz fehlt in der Bank (Platzhalter: Kettenname) – Nr. 1
%% TODO Anweisung: ein Satz über der ersten Teilaufgabe fehlt in der Bank – Nr. 1
\begin{aufgabe}{Kumulierte Binomialwahrscheinlichkeiten berechnen und lesen}
\begin{teile}
\teil $X$ ist binomialverteilt mit $n = 3$ und $p = 0{,}5$. $P(X \le 1)$?
\teil $X$ ist binomialverteilt mit $n = 20$ und $p = 0{,}3$. $P(X \le 4)$ auf vier Stellen?
\end{teile}
\end{aufgabe}

%% TODO Titel: Ich-kann-Satz fehlt in der Bank (Platzhalter: Kettenname) – Nr. 2
%% TODO Anweisung: ein Satz über der ersten Teilaufgabe fehlt in der Bank – Nr. 2
\begin{aufgabe}{Gegenereignis und Komplementärsummen („mindestens“ gegen „höchstens“)}
\begin{teile}
\teil $P(X \le 5) = 0{,}8$. $P(X \ge 6)$?
\teil $X$ ist binomialverteilt mit $n = 50$ und $p = 0{,}2$; $P(X \le 12) \approx 0{,}8139$. $P(X \ge 13)$?
\end{teile}
\end{aufgabe}

%% TODO Titel: Ich-kann-Satz fehlt in der Bank (Platzhalter: Kettenname) – Nr. 3
%% TODO Anweisung: ein Satz über der ersten Teilaufgabe fehlt in der Bank – Nr. 3
\begin{aufgabe}{Erwartungswert n · p als Orientierung, auf welcher Seite der Ablehnungsbereich liegt}
\begin{teile}
\teil $n = 60$, $p = 0{,}5$: Erwartungswert $n \cdot p$?
\teil $n = 900$, $p = 0{,}035$: Erwartungswert $n \cdot p$?
\end{teile}
\end{aufgabe}

%% TODO Titel: Ich-kann-Satz fehlt in der Bank (Platzhalter: Kettenname) – Nr. 4
%% TODO Anweisung: ein Satz über der ersten Teilaufgabe fehlt in der Bank – Nr. 4
\begin{aufgabe}{Funktionsgraphen ablesen und deuten (Werte, Monotonie)}
\begin{teile}
\teil Lies $f(2)$ ab.

\begin{ksys}[xmin=0,xmax=8,ymin=0,ymax=5,ablesen]\funktion{0.5*\x+1}{f}\end{ksys}
\teil Lies ab: Für welches $x$ ist $f(x) = 1$?

\begin{ksys}[xmin=0,xmax=8,ymin=0,ymax=5,ablesen]\funktion{4/(\x+1)}{f}\end{ksys}
\end{teile}
\end{aufgabe}

%% TODO Titel: Ich-kann-Satz fehlt in der Bank (Platzhalter: Kettenname) – Nr. 5
%% TODO Anweisung: ein Satz über der ersten Teilaufgabe fehlt in der Bank – Nr. 5
\begin{aufgabe}{Den Rechner für kumulierte Verteilungswahrscheinlichkeiten einsetzen}
\begin{teile}
\teil $X$ ist binomialverteilt mit $n = 10$ und $p = 0{,}5$. $P(X \le 3)$ mit dem Rechner?
\teil $X$ ist binomialverteilt mit $n = 120$ und $p = 0{,}35$. $P(35 \le X \le 45)$?
\end{teile}
\end{aufgabe}

%% TODO Titel: Ich-kann-Satz fehlt in der Bank (Platzhalter: Kettenname) – Nr. 6
%% TODO Anweisung: ein Satz über der ersten Teilaufgabe fehlt in der Bank – Nr. 6
\begin{aufgabe}{Gegenereignis und Komplementärsummen („mindestens“ gegen „höchstens“) – Fehler finden}
\begin{teile}
\teil $X$ ist binomialverteilt mit $n = 30$ und $p = 0{,}2$. Tom rechnet: $P(X \ge 10) = 1 - P(X \le 10) \approx 1 - 0{,}9744 = 0{,}0256$. Finde den Fehler und rechne richtig.
\end{teile}
\end{aufgabe}

%% TODO Titel: Ich-kann-Satz fehlt in der Bank (Platzhalter: Kettenname) – Nr. 7
%% TODO Anweisung: ein Satz über der ersten Teilaufgabe fehlt in der Bank – Nr. 7
\begin{aufgabe}{Gegenereignis und Komplementärsummen („mindestens“ gegen „höchstens“) – selbst rechnen}
\begin{teile}
\teil $X$ ist binomialverteilt mit $n = 40$ und $p = 0{,}25$; $P(X \le 11) \approx 0{,}7151$. $P(X \ge 12)$?
\end{teile}
\end{aufgabe}
